\documentclass[a4paper,12pt]{article}

\usepackage[T1]{fontenc}
\usepackage[italian]{babel}
\usepackage{graphicx}
\usepackage{geometry}
\usepackage{tabularx}
\usepackage[table]{xcolor}
\usepackage[most]{tcolorbox}
\usepackage{fancyhdr}
\usepackage[hidelinks]{hyperref}

\newcommand{\groupmail}{alittlebyte.swe@gmail.com}
\newcommand{\grouplogo}{../../../.github/site-src/assets/logo_alittlebyte.png}
\newcommand{\groupsymbol}{../../../.github/site-src/assets/solo_logo.png}

\geometry{
    left=2.5cm,
    right=2.5cm,
    top=2.5cm,
    bottom=2.5cm,
    headheight=331pt,
    includeheadfoot
}

\pagestyle{fancy}
\fancyhf{}

\renewcommand{\headrulewidth}{0pt}
\fancyhead{}

\renewcommand{\footrulewidth}{0.4pt}
\fancyfoot[L]{\includegraphics[height=0.6cm]{\groupsymbol}\hspace{0.45em}aLittleByte}
    \fancyfoot[C]{\thepage}
\fancyfoot[R]{Verbale 2026-06-09}

\fancypagestyle{firstpage}{
    \fancyhf{}
    \renewcommand{\headrulewidth}{0pt}
    \renewcommand{\footrulewidth}{0.4pt}
    \fancyhead[C]{%
        \makebox[\textwidth][c]{%
            \begin{tabular}{c}
                \includegraphics[height=10cm]{\grouplogo}\\[1ex]
                \small\texttt{\groupmail}
            \end{tabular}%
        }
    }
    \fancyfoot[L]{\includegraphics[height=0.6cm]{\groupsymbol}\hspace{0.45em}aLittleByte}
        \fancyfoot[C]{\thepage}
    \fancyfoot[R]{Verbale 2026-06-09}
}

\providecommand{\glslink}[2]{\href{https://alittlebyte-19.github.io/Documentazione/glossario.html\#gls-#1}{\underline{#2}\textsuperscript{\scriptsize G}}}

\begin{document}

\thispagestyle{firstpage}

\vspace*{0.5cm}

\begin{center}
    {\LARGE \textbf{Verbale Interno - 2026-06-09}}
\end{center}

\vspace{1.5cm}

\begin{center}
    \renewcommand{\arraystretch}{1.3}
    \begin{tcolorbox}[
        width=0.92\textwidth,
        colback=gray!6,
        colframe=gray!45,
        boxrule=0.5pt,
        arc=2mm,
        boxsep=0pt,
        left=8pt, right=8pt, top=8pt, bottom=8pt
    ]
        \begin{tabularx}{\linewidth}{>{\bfseries}p{3.5cm} X}
            Gruppo      & aLittleByte (gruppo 19) \\
            Redattore   & V. Y. Kaneda Fernandes \\
            Verificatori & Alex Oloni\\
            Destinatari & Prof. Tullio Vardanega, Prof. Riccardo Cardin \\
            Data        & 2026-06-09\\
        \end{tabularx}
    \end{tcolorbox}
\end{center}

\newpage

\newgeometry{left=2.5cm,right=2.5cm,top=2.5cm,bottom=2.5cm,includefoot}
\tableofcontents
\newpage
\section{Informazioni Generali}
\section*{Specifiche Documento}
\hrule
\vspace{1cm}

\begin{tabular}{>{\bfseries}p{5cm} | p{10cm}}
Luogo Riunione & \glslink{discord}{Discord} \\[6pt]

Orario (Inizio - Fine) & 15:00 -- 15:30 \\[6pt]

\glslink{redattore}{Redattore} &
  V. Y. Kaneda Fernandes\\[6pt]

\glslink{amministratore}{Amministratore} & 
A. Maule\\[6pt]

\glslink{verificatore}{Verificatore} &
A. Oloni\\[6pt]

Partecipanti &
  A. Oloni\\ &
  L. Soligo\\&
  A. Gastaldello\\&
  E. Bellet\\&
  D. Belluz\\&
  A. Maule\\&
  V. Y. Kaneda Fernandes\\[6pt]

\end{tabular}


\section{Ordine del Giorno}
    \begin{itemize}
    \item Analisi dell'operato dopo la valutazione insufficiente \\sull'\glslink{rtb-requirements-and-technology-baseline}{RTB (Requirements and Technology Baseline)}.
    \item Sistemazione dell' \glslink{analisi-dei-requisiti-adr}{Analisi dei Requisiti}
    \item Adattamenti del PoC
    \item Valutazione su un nuovo incontro con l'azienda
\end{itemize}


\section{Attività svolte}
\subsection{Analisi dell'operato dopo la valutazione insufficiente \\sull'\glslink{rtb-requirements-and-technology-baseline}{RTB (Requirements and Technology Baseline)}.}
Il team ha discusso principalmente sull'esito negativo avuto con l'incontro con il prof. Cardin, principalmente riguardanti alcune valutazioni tecnologie ritenute non sufficienti o con alternative migliori, oltre a varie sviste sull'AdR.

\subsection{Sistemazione dell' \glslink{analisi-dei-requisiti-adr}{Analisi dei Requisiti}}
Una delle critiche principali dopo la consegna sono state alcuni errori dentro il documento dell' \glslink{analisi-dei-requisiti-adr}{Analisi dei Requisiti}, sulla rappresentazione degli Use Case e su alcune svisste riguardo i requisiti. 
Si riporta quindi che andrà sistemato e verificato da tutti i componenti del gruppo, esprimendo e discutendo la best practice, per diminuire errori e che la rappresentazione sia chiara.

\subsection{Adattamenti del PoC}
Si pone quindi come obiettivo principale dell'Inizio dello sprint attuale, confermato in maniera unanime da tutti i membri del gruppo, di migliorare la struttura del Proof of Concept, sistemando e correggendo il codice attualmente
utilizzato e vengono valutate alternative su alcuni aspetti del PoC. 

\subsection{Valutazione su un nuovo incontro con l'azienda}
Viene discusso se chiedere nuove lucidazioni all'azienda proponente, affinchè le scelte che andremo a prendere non vadano ad intaccare sulla richiesta di quest'ultima. Si è deciso quindi di valutare prima le parti che andranno sistemate, e se ci
sono cambiamenti importanti, si consulterà l'azienda tramite riunione o i vari modi di comunicazione che sono stati definiti con codesta.

\section{To-Do List}
\begin{table}[h]
    \centering
    \begin{tabular}{|p{4.5cm}|p{5cm}|l|}
    \hline
        \rowcolor{lightgray}\textit{\textbf{Persona Incaricata}}  & \textbf{Task Assegnata} &\textbf{Link \glslink{issue}{Issue}} \\
    \hline
        \textit{A. Oloni e D. Belluz} & Sistemazione Analisi dei Requisiti& \href{https://github.com/aLittleByte-19/Documentazione/issues/121}{\#121}\\
    \hline 
        \textit{V. Y. K. Fernandes} & \glslink{diario-di-bordo}{Diario di Bordo} 12/06/2026 & \href{https://github.com/aLittleByte-19/Documentazione/issues/123}{\#123}\\
    \hline
        \textit{[Tutti i membri del gruppo]} & Correzione e sistemazione del PoC & \href{https://github.com/aLittleByte-19/Documentazione/issues/124}{\#124}\\
    \hline
        \textit{V. Y. K. Fernandes} & Stesura \glslink{verbale-interno}{Verbale Interno} 09/06/2026 & \href{https://github.com/aLittleByte-19/Documentazione/issues/122}{\#122}\\
    \hline
    \end{tabular}
    \label{tab:placeholder}
\end{table}

\end{document}